\documentclass[11pt]{article}
\usepackage[margin=1in]{geometry}
\usepackage{amsmath,amssymb}
\usepackage{setspace}
\usepackage{booktabs}
\usepackage[colorlinks=true, linkcolor=blue, citecolor=black, urlcolor=blue]{hyperref}
\usepackage{xcolor}

\onehalfspacing
\title{Who Pays for Senior Tax Exemptions?\\[0.5em]\large Brainstorming Document}
\author{Jeffrey Ohl}
\date{January 2026}

\begin{document}
\maketitle

\section*{The Idea}

Conway and Rork (2012) show state income tax exemptions for seniors don't attract migrants. We ask: who pays for them anyway?

If exemptions don't attract seniors but do cost revenue, the burden falls entirely on existing residents---through higher taxes or reduced services. The fiscal competition justification fails, but the fiscal cost remains.

\textbf{Variation}: Use Conway's policy coding (state-level changes in senior income tax treatment).

\textbf{New LHS}: Budget outcomes---tax rates on other groups, spending levels, budget composition.

\subsection*{What's the Contribution?}

After Conway, the reader's prior should already be ``exemptions cost money and someone pays.'' That's a logical implication of the null migration result. So the contribution isn't showing that \textit{someone} pays---it's showing \textit{who} pays and \textit{how much}.

The interesting findings would be:
\begin{itemize}
    \item \textbf{Which group loses}: Education? Low-income taxpayers (via sales tax)? Working-age income taxpayers?
    \item \textbf{Magnitude}: Is it dollar-for-dollar pass-through, or do states absorb part of the cost through other channels?
    \item \textbf{Heterogeneity}: Do states with different political compositions adjust differently?
\end{itemize}

A finding like ``senior income tax exemptions reduce K-12 spending by \$X per dollar of exemption'' is specific and policy-relevant. A finding like ``exemptions are offset by higher taxes somewhere'' is less interesting---we already expected that.

\section{Magnitude Check: Are We Powered?}

\subsection{How Big Are Senior Income Tax Exemptions?}

\textbf{What we think we know} (needs verification):
\begin{itemize}
    \item 41 states exempt Social Security benefits entirely
    \item $\sim$7\% of state income tax revenue goes to senior exemptions (source?)
    \item Some states also exempt pension income, IRA distributions, etc.
\end{itemize}

\textbf{BOTEC}:
\begin{itemize}
    \item Total state income tax revenue (2023): $\sim$\$500 billion
    \item If 7\% goes to senior exemptions: $\sim$\$35 billion nationally
    \item Average state: $\sim$\$700 million/year
    \item But highly variable---states with no income tax contribute zero; high-tax states contribute more
\end{itemize}

\textbf{As share of state budgets}:
\begin{itemize}
    \item Average state general fund: $\sim$\$20--30 billion
    \item \$700M / \$25B = 2.8\% of general fund
    \item Not trivial, but not 10\%
\end{itemize}

\textcolor{red}{\textbf{TODO}: Verify these numbers. Find actual estimates of senior income tax expenditures by state.}

\subsection{Power Concerns}

If senior exemptions are 2--3\% of state budgets, can we detect offsetting changes in taxes or spending?

\textbf{Optimistic case}: A state expanding senior exemptions by 1pp of budget should show up as either:
\begin{itemize}
    \item 1pp higher taxes elsewhere, or
    \item 1pp lower spending somewhere, or
    \item Some combination
\end{itemize}

\textbf{Pessimistic case}: State budgets are noisy. Many things change simultaneously. A 1--2\% shift may be hard to detect, especially with limited policy variation (how many states changed senior exemptions in Conway's sample?).

\textbf{What helps}:
\begin{itemize}
    \item Panel data over multiple decades (more observations)
    \item Sharp policy changes (discrete events, not gradual drift)
    \item Pre-trends analysis (did budgets diverge before the policy change?)
\end{itemize}

\section{Framing Different Findings}

\subsection{Finding: Self-Financing (Null on Budget)}

\textit{Senior exemptions don't change taxes or spending levels.}

\textbf{Interpretation}: States absorb the cost through:
\begin{itemize}
    \item Revenue growth (economy expands, new revenues offset exemption cost)
    \item Efficiency gains (cut waste rather than services)
    \item Deficit financing / rainy day fund drawdowns
\end{itemize}

\textbf{Policy implication}: Exemptions are ``free'' in equilibrium---no visible losers.

\textbf{Problem with this finding}: Hard to publish a null. Also, ``absorbed through revenue growth'' is unsatisfying---the counterfactual is that revenue growth could have funded something else.

\subsection{Finding: Education Spending Falls}

\textit{States that expand senior exemptions cut education spending.}

\textbf{Interpretation}: Intergenerational conflict. Seniors vote for tax breaks for themselves, financed by reduced investment in children.

\textbf{Why education?}
\begin{itemize}
    \item Largest discretionary category in most state budgets
    \item Seniors don't directly benefit
    \item Politically easier to cut than Medicaid (federal matching) or public safety
\end{itemize}

\textbf{Policy implication}: Senior exemptions have hidden costs borne by future generations.

\textbf{Literature connection}: ``Gray peril'' hypothesis---aging populations underinvest in education. This would be causal evidence.

\subsection{Finding: Other Taxes Rise}

\textit{States that expand senior income tax exemptions raise property taxes, sales taxes, or income taxes on non-seniors.}

\textbf{Interpretation}: Tax shifting. The exemption doesn't reduce total tax burden; it shifts who pays.

\textbf{Which taxes?}
\begin{itemize}
    \item Property taxes: States could reduce aid to localities, forcing property tax increases
    \item Sales taxes: Regressive; burden falls on low-income households
    \item Income taxes on working-age: Direct substitution within same tax base
\end{itemize}

\textbf{Policy implication}: Senior exemptions are a transfer from [X] to seniors, not a tax cut.

\subsection{Finding: Nothing Detectable (True Null)}

\textit{No measurable effect on any budget category.}

\textbf{Interpretation options}:
\begin{itemize}
    \item Effects too small / diffuse to detect (power problem)
    \item Effects spread across many categories (no single category moves enough)
    \item States adjust on margins we don't observe (hiring freezes, deferred maintenance)
\end{itemize}

\textbf{What to do}: Report confidence intervals. ``We can rule out effects larger than X\% of spending.''

\section{Identification Concerns}

\textbf{Endogeneity}: States don't randomly adopt senior exemptions. They adopt when:
\begin{itemize}
    \item Senior population is large/growing (political pressure)
    \item Budgets are flush (can afford it)
    \item Competing with neighbors for retirees
\end{itemize}

These confounders correlate with budget outcomes directly.

\textbf{Conway's approach}: State fixed effects + time fixed effects. Identifies off within-state changes over time. But still vulnerable to time-varying confounders (states adopt exemptions during booms, then cut spending during busts---but the bust caused both).

\textbf{Possible improvements}:
\begin{itemize}
    \item Event study with pre-trends (did spending already diverge before adoption?)
    \item Instrument? (Hard to think of one)
    \item Synthetic control for individual large policy changes
\end{itemize}

\section{Data Requirements}

\textbf{Policy variation}: Need Conway's coding of state senior income tax treatment, or reconstruct it.
\begin{itemize}
    \item Social Security exemption (full, partial, none)
    \item Pension income exemption
    \item Other senior-specific provisions
    \item Timing of changes
\end{itemize}

\textbf{Budget data}: Census of Governments / Annual Survey of State and Local Government Finances
\begin{itemize}
    \item State tax revenue by source (income, sales, property)
    \item State spending by category (education, health, infrastructure)
    \item State aid to local governments
\end{itemize}

\textbf{Controls}:
\begin{itemize}
    \item Demographics (age structure, income)
    \item Economic conditions (unemployment, GDP growth)
    \item Political variables (governor party, legislature composition)
\end{itemize}

\section{Next Steps}

\begin{enumerate}
    \item \textbf{Verify magnitudes}: Find actual estimates of senior income tax expenditures by state. Is this 2\% of budgets or 5\%?
    \item \textbf{Get Conway's data}: Check if policy coding is available in appendix/replication files. If not, reconstruct from Tax Foundation or CCH state tax guides.
    \item \textbf{Count policy changes}: How many states changed senior income tax treatment in the relevant period? Is there enough variation?
    \item \textbf{Preliminary budget analysis}: Do simple pre/post comparisons for a few large policy changes to see if anything moves.
\end{enumerate}

\section{Honest Assessment}

This is a harder identification problem than the Florida capitalization question. State budgets are noisy, policy changes are endogenous, and the magnitude may be too small to detect.

But the question is more interesting: ``Who pays for senior tax breaks?'' has broader policy relevance than ``Do property tax exemptions capitalize?''

The Conway null on migration is the setup. If exemptions don't attract seniors, they're pure transfers---and documenting who loses is valuable.

\end{document}
